Out of the \data{audit/sample_count_int} LLM-validated transcripts with the highest scores, \data{audit/false_positive_count_int} were classified by the human auditor as false positives. Of those, 29 lacked any relevant content and have been interpreted as hallucinations. The other 163 contain some content which, in principle, could be part of a collusive message. Specifically, the transcripts refer to the firm's competitive conduct or the state of market competition but did not connect it to the conduct of competitors or the industry and thereby fell short of facilitating coordinated conduct. Here we describe the most common content that appeared in the LLM's excerpts and seemed to be the reason for the LLM assigning a high score.

67 transcripts referred to a firm reducing, restricting, or managing its supply or capacity in a manner consistent with higher prices or avoiding lower prices. However, it did not say or suggest that this conduct is something the industry should or will or needs to do.

\begin{quote}
``Our intention [is] to maintain a disciplined sales approach. Rather than selling at material decrease prices. We implemented a strategy to hold back product [and] resulting in a significant buildup of the inventories.'' \\
\mbox{[\data{quoted_excerpt_labels/audit_errors/q298141}]}
\end{quote}

\begin{quote}
``We will continue to manage our customer needs and will not restart capacity on a whim just to add tonnage to the spot market. That would not be good for anyone. Value over volume is a simple philosophy that we will carry on into the future, and it's why you are not going to hear me talk about capacity utilization or even market share.'' \\
\mbox{[\data{quoted_excerpt_labels/audit_errors/q445982}]}
\end{quote}

\begin{quote}
``We stand by our strategy of not forsaking value for volume and protecting our profitability with continued actions on our cost structure [and we] cut our production capacity.'' \\
\mbox{[\data{quoted_excerpt_labels/audit_errors/q458895}]}
\end{quote}

\begin{quote}
``We go down by elevator and we go up by stairs. So to avoid this huge volatility, we are aiming … to reduce capacity.'' \\
\mbox{[\data{quoted_excerpt_labels/audit_errors/q198201}]}
\end{quote}

By noting its leadership position or forecasting constrained industry supply, some announcements do get close to suggesting that firms coordinate on reducing supply or capacity.

\begin{quote}
``Our first step is to stop the significant price declines. The most important tool that we have to accomplish this is to reduce capacity to get glass supply in line with market demand. When you're more than half the market, that really adds up.'' \\
\mbox{[\data{quoted_excerpt_labels/audit_errors/q252936}]}
\end{quote}

\begin{quote}
``And until now, we have not heard anyone else announcing any capacity enhancements in the short term. None of them is increasing capacity.'' \\
\mbox{[\data{quoted_excerpt_labels/audit_errors/q452530}]}
\end{quote}

\begin{quote}
``Additional capacity reductions can be expected [and] we forecast an extended period when the … industry operates a smaller, more profitable footprint. Our outlook is for lumber pricing to gradually increase.'' \\
\mbox{[\data{quoted_excerpt_labels/audit_errors/q654876}]}
\end{quote}

48 of the transcripts refer to higher prices - whether noting having raised prices or expressing a need to raise prices - or to pricing power. Such a statement might be made in the context of reducing supply.

\begin{quote}
``We reduced our … production capacity by more than 20\% [and] we have already announced another 20\% price hike. We have started to see the sign of improving pricing trend.'' \\
\mbox{[\data{quoted_excerpt_labels/audit_errors/q638665}]}
\end{quote}

\begin{quote}
``We're going to... be price stewards. We don't chase volume with price.'' \\
\mbox{[\data{quoted_excerpt_labels/audit_errors/q422806}]}
\end{quote}

\begin{quote}
``We have taken a call … to increase prices of cement to help manage the rising cost of production.'' \\
\mbox{[\data{quoted_excerpt_labels/audit_errors/q320434}]}
\end{quote}

With a reference to competitors reducing capacity in the context of an emphasis on raising prices, this announcement is getting close to having collusive content.

\begin{quote}
``In the short term, there's only one priority, and that's price increases and price increase and price increases. We have good capacity utilization. We have strong brands. We always go for pricing every time. Many competitors, especially in the tissue area, announced that they were shutting down paper machines due to this is looked at an Italian player ... and I think we'll see more of that.'' \\
\mbox{[\data{quoted_excerpt_labels/audit_errors/q215021}]}
\end{quote}

In several transcripts, a firm announced it was striving to avoid lowering prices.

\begin{quote}
``We concluded that lowering price to gain share in a market like we are experiencing this year was not the right strategy.'' \\
\mbox{[\data{quoted_excerpt_labels/audit_errors/q375619}]}
\end{quote}

\begin{quote}
``We don't want to underprice and overshoot on demand. We keep some of the exclusive capacity in order to be able to sell at high prices. We want to be profitable [even] if we lost a couple of market share percentage points.'' \\
\mbox{[\data{quoted_excerpt_labels/audit_errors/q257163}]}
\end{quote}

\begin{quote}
``[We are] idling significant chunks of our assets and slashing our participation in weak markets. We're not just going to participate in those weak markets because it's a chase down into the mud.'' \\
\mbox{[\data{quoted_excerpt_labels/audit_errors/q662160}]}
\end{quote}

A firm indicated a strategy to strengthen pricing power.

\begin{quote}
``We deliberately allocate inventory to fiscal years and regions to ensure we maintain apparent scarcity and, therefore, protect our brand health and protect our pricing power.'' \\
\mbox{[\data{quoted_excerpt_labels/audit_errors/q631021}]}
\end{quote}

\begin{quote}
``We're in the early stages of having a terrific opportunity over the next 5 years to be in a position where we can increase the returns at our disposal facilities, primarily through price. We’re one of the only players with a little bit of excess capacity left.'' \\
\mbox{[\data{quoted_excerpt_labels/audit_errors/q190678}]}
\end{quote}

A firm described how it is being more disciplined and less aggressive which, if it had expressed competitors acting or needing to act in a similar manner, would have encompassed collusive content.

\begin{quote}
``We do not want to accelerate a rapid decline in price of the products. We will only ramp it up very moderately when we start production because we do not want to accelerate a collapse in price. We want to be more disciplined in our approach.'' \\
\mbox{[\data{quoted_excerpt_labels/audit_errors/q604617}]}
\end{quote}

\begin{quote}
``We are not chasing market share. We are focusing on quality customers and then driving spend on our network by our existing customers.'' \\
\mbox{[\data{quoted_excerpt_labels/audit_errors/q643679}]}
\end{quote}

\begin{quote}
``We have … implemented Microchip's disciplined pricing process … to ensure that we are competitive and don't chase bad business in the pursuit of profitless prosperity.'' \\
\mbox{[\data{quoted_excerpt_labels/audit_errors/q320228}]}
\end{quote}

In this case, the firm is expressing an aspiration for the industry which is getting close to collusive content.

\begin{quote}
``We would like to drive the industry in a direction where we can move back to more positive margins.'' \\
\mbox{[\data{quoted_excerpt_labels/audit_errors/q550676}]}
\end{quote}

A firm noted that it has market power or what the sources of market power are. While in one case it does mention reduced industry capacity being a source, it is of the form of a factual observation.

\begin{quote}
``Supply is having a harder time to react to demand [and] that's [where] we're getting the pricing power.'' \\
\mbox{[\data{quoted_excerpt_labels/audit_errors/q401404}]}
\end{quote}

\begin{quote}
``When our actual sales pace exceeds our budgeted sales pace, we try to push prices.'' \\
\mbox{[\data{quoted_excerpt_labels/audit_errors/q651762}]}
\end{quote}

\begin{quote}
``We expect to continue to implement favorable pricing actions … due to capacity constraints in our industry. We are operating today in a very strong demand environment without enough industry capacity to meet that demand and expect that to continue into the future.'' \\
\mbox{[\data{quoted_excerpt_labels/audit_errors/q593253}]}
\end{quote}

Some transcripts refer to reduced industry supply or capacity but only in the context of stating it as a fact. Still, the announcements are getting close to content that may encourage competitors to reduce supply or capacity.

\begin{quote}
``This moment of rationalization of supply is new in the history of the Brazilian civil aviation industry. We wouldn't like to start or restart price war in the Brazilian market. Our strategy of reducing capacity is indeed in place and it's showing results.'' \\
\mbox{[\data{quoted_excerpt_labels/audit_errors/q481706}]}
\end{quote}

\begin{quote}
``[There is] limited new barge construction in the industry and many units going in for maintenance. So we should have a net decline this year in supply. Nobody is really panicking and trying to go out and build any new equipment right now.'' \\
\mbox{[\data{quoted_excerpt_labels/audit_errors/q327107}]}
\end{quote}

In 14 transcripts, the firm refers to industry consolidation by noting that it has been consolidating, needs to consolidate, or will consolidate. In some instances, the firm attributes consolidation to too much supply or capacity or too many firms. In one case, they do seem to be facilitating such consolidation.

\begin{quote}
``We expect the industry will consolidate to approximately 5 Pan-India players. Once that happens, the industry will further see tariff hardening, with pricing power returning.'' \\
\mbox{[\data{quoted_excerpt_labels/audit_errors/q321900}]}
\end{quote}

\begin{quote}
``In order to really improve value creation, we need consolidation and further restructuring. There’s a need for further consolidation and restructuring in this industry to be able to really actually take the costs down and balance the markets as well.'' \\
\mbox{[\data{quoted_excerpt_labels/audit_errors/q414209}]}
\end{quote}

In summing up, the LLM identified instances in which a firm’s message was consistent with less competition in that it referred to reducing supply or capacity or raising or at least not lowering prices. To facilitate collusive conduct, a firm needs to link such content to industry conduct in a manner that encourages similar conduct by competitors. The LLM failed in not verifying that such a link was present in the firm’s message. Consequently, the LLM had a high false positive rate.
